\section{Discussion}
\label{sec:discussion}
% =====================================================================

\paragraph{What is robust.} The comparison was repeated with three derivative estimators, two covariances, four fiducial HODs, four or six parameters, two tracer populations and several data vectors. Two conclusions never changed: ASTRA split and density split carry comparable information (error ratios 0.75--1.15 in the DESI-like analysis with the independent covariance, 0.6--2.1 in the reliable Quijote forecasts), and they are complementary. The ranking for individual parameters changed with the analysis choices by 10--40\% (Table~\ref{tab:evolution}), the same order as the differences between the methods, so we do not claim an overall winner. Two trends are consistent across analyses: density split is slightly ahead on $\omega_b$ and $n_s$ (and $\Omega_m$ in Quijote) with the random-quartile cross-correlations, and ASTRA is ahead on $\sigma_8$ through the auto-correlations and dense-region environments.

\paragraph{What we learnt about the methodology.} (i) Emulator suites in which the HOD is drawn independently in each cosmology do not allow naive finite differences: the HOD mismatch dominates the cosmological response by an order of magnitude, and paired differences do not cancel the galaxy sampling noise of different HOD catalogues. (ii) Single-box covariances, even with many subvolumes and lines of sight, can bias the \emph{comparison} of statistics, not only their absolute errors. (iii) Environment statistics of sparse tracers have noisy derivatives, which a six-parameter Fisher matrix amplifies along its degenerate directions; comparing two Fisher estimators that bracket the truth identifies the reliable forecasts. (iv) The clustering statistics respond essentially to a single combination of the 12 HOD parameters.

\paragraph{Limitations.} (i) The forecasts are for one $(500\Mpch)^3$ volume, where the errors are comparable to the prior widths: they compare information content and are not realistic constraints. (ii) The inference is calibrated but not yet discriminating. (iii) The Quijote halos are ten times sparser than the LRGs, and ASTRA's adaptive scale is correspondingly larger. (iv) The emulator suite has a single HOD realisation per point, whose galaxy sampling noise sets the floor of the emulator error.

\paragraph{Improvements.} Several HOD realisations (seeds) per point, or HOD-matched catalogues generated with AbacusHOD in all cosmologies, would lower the emulator error, making an inference on a DESI-like volume discriminating, and give clean derivatives in the DESI-like setting. Larger sets of independent boxes would remove the remaining dependence on the covariance.

\paragraph{Beyond quartiles.} Being local and free of smoothing, ASTRA provides more than a density ranking: a classification of every galaxy into voids, sheets, filaments and knots, a continuous environment value usable as a mark, and a natural treatment of survey masks through the random points. The test of the classes on Quijote suggests that this is worth pursuing; natural next steps are an equal-size two-way split to understand their gain, the probabilistic classes of \cite{ASTRA} (classification repeated over many random catalogues), and an emulator of class statistics for DESI.


